% Generated draft: requires mathematical and source review.
\documentclass[11pt]{article}
\usepackage[margin=25mm]{geometry}
\usepackage{fontspec}
\setmainfont{Latin Modern Roman}
\usepackage{amsmath,amssymb,mathtools,unicode-math}
\setmathfont{Latin Modern Math}
\usepackage{enumitem,needspace,xurl,hyperref}
\setlength{\parindent}{0pt}
\setlength{\parskip}{5pt}
\setlength{\emergencystretch}{3em}
\newcommand{\N}{\mathbb{N}}
\newcommand{\Z}{\mathbb{Z}}
\newcommand{\Q}{\mathbb{Q}}
\newcommand{\R}{\mathbb{R}}
\newcommand{\C}{\mathbb{C}}
\newcommand{\F}{\mathbb{F}}
\newcommand{\A}{\mathbb{A}}
\newcommand{\PP}{\mathbb{P}}
\newcommand{\E}{\mathbb{E}}
\newcommand{\eps}{\varepsilon}
\newcommand{\dd}{\,\mathrm{d}}
\begin{document}
\section*{UnsolvedMath ID: 30001600}
\textbf{Source problem code:} catalogue entry
\subsection{Definitions and mathematical statement}
The supplied record concerns a smooth curve \(C\), a line bundle \(L\) on \(C\), its gonality \(\delta\), and the integer
\[
V=H^0(C,L),\qquad r=\dim V-1.
\]
The ground field, its characteristic, and the projectivity and geometric connectedness assumptions on \(C\) are not specified. Consequently, the following gives a conditional formulation, rather than silently supplying these missing hypotheses.

If \(C\) is interpreted as a smooth projective geometrically integral curve over a field \(k\), and gonality is interpreted over \(k\), define
\[
\delta=\min\{\deg(f):f\colon C\longrightarrow\mathbb P^1_k\text{ is a nonconstant }k\text{-morphism}\}.
\]
Whether the intended gonality is instead computed after extending to an algebraic closure remains unresolved. All tensor products and exterior powers below are over the chosen ground field.

For a positive integer \(p\), the Koszul cohomology group \(K_{p,1}(C,L)\) is the homology at the middle term of the complex
\[
\bigwedge^{p+1}V\otimes H^0(C,\mathcal O_C)
\xrightarrow{d}
\bigwedge^p V\otimes H^0(C,L)
\xrightarrow{d}
\bigwedge^{p-1}V\otimes H^0(C,L^{\otimes 2}),
\]
where \(d\) is the alternating Koszul differential induced by multiplication of sections. Thus \(K_{p,1}(C,L)\neq 0\) means that a nonzero linear syzygy occurs at homological index \(p\).

The question attributed in the record to Green--Lazarsfeld asks whether, for a sufficiently positive \(L\),
\[
\forall p\in\mathbb Z_{\geq 1},\qquad
K_{p,1}(C,L)\neq 0
\quad\Longleftrightarrow\quad
1\leq p\leq r-1-\delta.
\]
The upper endpoint \(r-1-\delta\) is retained exactly as supplied, with \(r=h^0(C,L)-1\). Its compatibility with the intended Koszul indexing requires source verification; this draft makes no correction. If the upper endpoint is less than \(1\), the displayed interval is empty.

If ``sufficiently positive'' means sufficiently large degree for each fixed curve, the intended quantifiers would be: for every curve \(C\) in the intended class, there exists an integer \(d_0(C)\) such that every line bundle \(L\) on \(C\) with \(\deg L\geq d_0(C)\) satisfies the displayed equivalence. This degree-threshold interpretation, the intended class of curves, and any additional hypotheses on \(L\) are not specified in the supplied record and remain qualifications to this conditional formulation.
\subsection{Short English statement}
For a sufficiently positive line bundle on a smooth curve of gonality \(\delta\), the supplied question asks whether nonzero linear Koszul syzygies occur exactly at indices \(1\leq p\leq r-1-\delta\), where \(r=h^0(C,L)-1\). The endpoint is preserved verbatim. The field, curve hypotheses, gonality convention, positivity condition, and indexing require human review.
\subsection{Sources}
\begin{itemize}
\item Local registry: \texttt{data/unsolvedmath/problems.json}, supplied record with permanent UnsolvedMath ID \(30001600\), title ``Green--Lazarsfeld Gonality Conjecture,'' exact source problem code \texttt{catalogue entry}. The record attributes the conjecture to Green--Lazarsfeld; its separate proposer field is empty.
\item Supplied source citation: \emph{Mini-Workshop: Linear Series on Algebraic Varieties (2011)}, Oberwolfach Reports, pp.~2613--2650, DOI \texttt{10.4171/owr/2010/45}. \url{https://doi.org/10.4171/owr/2010/45}. The supplied background describes Oberwolfach Reports as documenting mathematical workshops and their open-problem discussions.
\item Supplied literature source: \url{https://arxiv.org/abs/math/0203246}. No title or author metadata for this reference is supplied.
\item Supplied literature source: \url{https://doi.org/10.1016/S1631-073X(03)00062-1}. No title or author metadata for this reference is supplied.
\item Literature assessment dated exactly \texttt{2026-08-21}: the supplied curation labels the problem ``partially solved'' and states that Green and Green--Lazarsfeld syzygy conjectures have been proved in broad generic and large-gonality regimes, while the extracted general statement is not fully resolved. This is supplied dated curation, not an independently verified literature assessment. The background explicitly requests independent verification of time-sensitive status and cited claims.
\item The registry labels the statement \texttt{corrected\_verified} and reports a manual semantic, typographic, and LaTeX review asserting that its cleanup preserves meaning and adds no unsupported source-dependent content. These are supplied review claims, not mathematical verification by this draft.
\end{itemize}
\end{document}
